\chapter{Lineární rovnice a stabilita}
\begin{chapquote}{Steven H. Strogatz}%The Joy Of X: A Guided Tour of Math, from One to Infinity
 „Velký, složitý lineární problém lze vždy rozdělit na menší, lépe zvládnutelné části. Každou část pak lze vyřešit samostatně a všechny dílčí výsledky lze následně zkombinovat, aby se vyřešil celý problém. U lineárního problému se celek přesně rovná součtu jednotlivých částí.“
\end{chapquote}
V této kapitole si ukážeme, že řešení jistého typu diferenciálních rovnic má jednoduchou geometrickou strukturu: stukturu \emph{vektorového prostoru}.
\section{Lineární homogenní rovnice}
\section{Lineární nehomogenní rovnice}
\begin{příklad}[Newtonův zákon ochlazování]
\begin{equation*}
	T'=-k(T-T_{out})
\end{equation*}
\end{příklad}
\begin{příklad}[Rezonance]
\begin{equation*}
  x''+\omega_0^{2}x=A\sin(\tilde{\omega}t)
\end{equation*}
\end{příklad}
\section{Soustavy lineárních rovnic}
